\newglossaryentry{aaa}{name={AAA},description={Producción de videojuego desarrollada con presupuestos elevados, equipos grandes y altos niveles de inversión técnica y artística. Se utiliza en la memoria para contrastar con el desarrollo independiente.}}
\newglossaryentry{almas}{name={Almas},description={Recurso principal de Via Inferni utilizado por el jugador para comprar objetos y activar ciertas acciones, como la recuperación de vida.}}
\newglossaryentry{algoritmodrunkardswalk}{name={Algoritmo Drunkard’s Walk},description={Algoritmo de generación procedural basado en recorridos aleatorios utilizado para crear mapas o distribuciones de salas conectadas. En el proyecto se emplea para garantizar conectividad y diversidad en los niveles.}}
\newglossaryentry{combateentiemporeal}{name={Combate en tiempo real},description={Sistema de combate donde las decisiones del jugador deben tomarse durante la acción sin pausas estratégicas.}}
\newglossaryentry{circulo}{name={Círculo},description={Nivel temático o bloque de progresión inspirado en la estructura narrativa del infierno.}}
\newglossaryentry{dualdadpersonajes}{name={Dualidad de personajes},description={Mecánica jugable basada en la alternancia entre Dante y Virgilio, cada uno con habilidades y funciones diferenciadas.}}
\newglossaryentry{dungeon}{name={Dungeon},description={Mazmorra o conjunto de salas conectadas utilizado como espacio principal de exploración en muchos videojuegos roguelike.}}
\newglossaryentry{enemigo}{name={Enemigo},description={Entidad hostil controlada por el juego que complica el avance del jugador y actúa como obstáculo jugable.}}
\newglossaryentry{feature}{name={Feature},description={Unidad funcional concreta dentro del desarrollo iterativo del proyecto. Cada feature representa una funcionalidad implementable y verificable.}}
\newglossaryentry{fdd}{name={Feature Driven Development (FDD)},description={Metodología de desarrollo basada en la implementación incremental de funcionalidades pequeñas y verificables.}}
\newglossaryentry{generacionprocedural}{name={Generación procedural},description={Técnica mediante la cual contenido del videojuego se genera algorítmicamente de forma automática y variable, permitiendo crear escenarios diferentes en cada partida.}}
\newglossaryentry{gitflow}{name={GitFlow},description={Modelo de organización de ramas en Git basado en ramas principales y ramas temporales para funcionalidades y correcciones.}}
\newglossaryentry{hud}{name={Head-Up Display (HUD)},description={Conjunto de elementos visuales que muestran información relevante del juego al jugador durante la partida, como vida, recursos o minimapa.}}
\newglossaryentry{indiegame}{name={Indie Game},description={Videojuego desarrollado por equipos pequeños o independientes, normalmente con recursos limitados.}}
\newglossaryentry{instanciacionsalas}{name={Instanciación de salas},description={Proceso mediante el cual las salas del mapa se generan y colocan dinámicamente durante la ejecución del juego.}}
\newglossaryentry{inventario}{name={Inventario},description={Sistema para almacenar y gestionar objetos recogidos durante la partida.}}
\newglossaryentry{llavecirculo}{name={Llave de círculo},description={Elemento de progresión que habilita el descenso al siguiente círculo una vez recogido dentro del nivel actual.}}
\newglossaryentry{maquinaestadosfinitos}{name={Máquina de estados finitos},description={Modelo de comportamiento basado en estados y transiciones utilizado para controlar la lógica de enemigos y entidades.}}
\newglossaryentry{menu}{name={Menú},description={Pantalla de navegación que permite acceder a opciones, pausas, configuración o acciones principales del juego.}}
\newglossaryentry{minimapa}{name={Minimapa},description={Representación simplificada del mapa del juego que facilita orientación y navegación al jugador.}}
\newglossaryentry{motorgrafico}{name={Motor gráfico},description={Software base utilizado para desarrollar videojuegos proporcionando renderizado, físicas, interfaz y gestión de escenas.}}
\newglossaryentry{objeto}{name={Objeto},description={Elemento interactivo o consumible que aporta un efecto concreto al jugador o a su progreso.}}
\newglossaryentry{observer}{name={Observer},description={Patrón de diseño basado en comunicación mediante eventos entre objetos desacoplados.}}
\newglossaryentry{permadeath}{name={Permadeath},description={Mecánica típica de los roguelike en la que la muerte implica reiniciar el progreso de la partida.}}
\newglossaryentry{pixelart}{name={Pixel Art},description={Estilo artístico basado en gráficos construidos mediante píxeles visibles, habitual en videojuegos retro y 2D.}}
\newglossaryentry{playtesting}{name={Playtesting},description={Proceso de prueba del videojuego mediante sesiones de juego orientadas a detectar errores y balancear la experiencia.}}
\newglossaryentry{pmbok}{name={PMBOK},description={Guía de referencia para la dirección de proyectos utilizada en la memoria como marco para seleccionar y adaptar procesos de planificación, seguimiento y control.}}
\newglossaryentry{prefab}{name={Prefab},description={Objeto reutilizable de Unity que almacena una configuración predefinida de componentes y elementos visuales, permitiendo instanciar entidades repetidamente dentro del juego.}}
\newglossaryentry{prototipado}{name={Prototipado},description={Proceso de construcción rápida de versiones funcionales preliminares para validar ideas y mecánicas.}}
\newglossaryentry{rejugabilidad}{name={Rejugabilidad},description={Capacidad de un videojuego para ofrecer experiencias diferentes y mantener interés tras múltiples partidas.}}
\newglossaryentry{roguelike}{name={Roguelike},description={Subgénero de videojuegos caracterizado por generación procedural, muerte permanente y alta dificultad.}}
\newglossaryentry{run}{name={Run},description={Partida individual completa dentro de un roguelike.}}
\newglossaryentry{sala}{name={Sala},description={Unidad espacial del nivel que agrupa un tipo concreto de interacción, reto o transición dentro de la partida.}}
\newglossaryentry{salaespecial}{name={Sala especial},description={Tipo de sala con función diferenciada dentro del mapa, como tienda, tesoro o sala secreta, que introduce recompensas, compras o exploración adicional.}}
\newglossaryentry{salajefe}{name={Sala del jefe},description={Sala asociada al tramo final de un círculo y al acceso que permite descender al siguiente nivel de la run.}}
\newglossaryentry{scriptableobject}{name={ScriptableObject},description={Tipo de recurso de Unity utilizado para almacenar datos compartidos desacoplados de la lógica de ejecución.}}
\newglossaryentry{sprite}{name={Sprite},description={Recurso gráfico 2D utilizado para representar personajes, enemigos, objetos o elementos de interfaz dentro del videojuego.}}
\newglossaryentry{tileset}{name={Tileset},description={Conjunto de piezas gráficas reutilizables empleadas para construir visualmente salas, suelos, muros y otros elementos del escenario.}}
\newglossaryentry{trunkbaseddevelopment}{name={Trunk-Based Development},description={Estrategia de integración frecuente basada en integrar cambios continuos sobre una rama principal compartida.}}
\newglossaryentry{uml}{name={UML},description={Lenguaje de modelado utilizado para representar gráficamente la arquitectura y relaciones entre componentes software.}}
\newglossaryentry{unity}{name={Unity},description={Motor de desarrollo de videojuegos utilizado como base técnica para implementar Via Inferni.}}
